% Irodov, Arihant PDF p. 30, printed p. 21.
% Start of section 1.3: work, kinetic energy, momentum.
\exproblem{118}{21}
质点沿 $xy$ 平面中的任意轨迹从
$\bm r_1=\bm i+2\bm j$ 移动到
$\bm r_2=2\bm i-3\bm j$，
其间受一恒力 $\bm F=3\bm i+4\bm j$。
上述各量均采用 SI 单位。求该恒力做的功。
\exsolution
恒力的功不依赖所走轨迹，仅由始末位移决定：
\[
\Delta\bm r=\bm r_2-\bm r_1
=(2-1)\bm i+(-3-2)\bm j
=\bm i-5\bm j.
\]
故
\[
\boxed{W=\int_1^2\bm F\cdot\mathrm d\bm r
=\bm F\cdot(\bm r_2-\bm r_1)
=3\times1+4\times(-5)=-17\,\mathrm J}.
\]
\excheck 负功表示该外力沿净位移方向的投影为负，
并不意味着质点不能按所给轨迹运动；题目还允许其他力同时作用。

\exproblem{119}{21}
一辆质量为 $m$ 的机车从静止开始运动，
其速率满足 $v=a\sqrt s$，
$s$ 为从起点起已经走过的路程，$a$ 为常数。
求开始运动后前 $t$ 秒内全部作用力做的总功。
\exsolution
注意本版 PDF 的文本提取可能漏掉 $\sqrt s$ 的根号；
原书第21印刷页的公式经图像核验确为 $v=a\sqrt s$。
沿单向运动有
\[
\frac{\mathrm ds}{\mathrm dt}=a\sqrt s.
\]
取“自 $t=0$ 即开始运动”的非静止解，
分离变量得
$2\sqrt s=at$，于是
\[
s(t)=\frac{a^2t^2}{4},\qquad
v(t)=a\sqrt{s(t)}=\frac{a^2t}{2}.
\]
由动能定理，所有力做功之和为动能增量：
\[
\boxed{W_{\rm total}(0,t)=\frac12mv(t)^2
=\frac{ma^4t^2}{8}}.
\]
\excheck $v=a\sqrt{s}$ 与 $v(0)=0$ 单独作为微分方程
仍有“先静止、再启动”的非唯一数学解；
物理上题目“开始运动后前 $t$ 秒”已选取立即起动的分支。

\exproblem{120}{21}
质量为 $m$ 的质点沿半径为 $R$ 的圆周运动，
其动能与走过的路程 $s$ 的关系为
$K(s)=as^2$，其中 $a>0$。
求合外力的大小作为 $s$ 的函数。
\exsolution
沿运动切线，合力切向分量由功—能关系确定：
\[
F_\tau=\frac{\mathrm dK}{\mathrm ds}=2as.
\]
又 $K=\frac12mv^2$，得 $mv^2=2as^2$，
故沿圆周法线方向的合力为
\[
F_n=\frac{mv^2}{R}=\frac{2as^2}{R}.
\]
两个分量互相垂直，因而合力大小
\[
\boxed{F(s)=\sqrt{F_\tau^2+F_n^2}
=2as\sqrt{1+\frac{s^2}{R^2}}}.
\]
\excheck 不能仅把 $F_\tau=\mathrm dK/\mathrm ds$ 当作合力；
曲线运动存在向心力分量。$a$ 的量纲为
$\mathrm{J/m^2}$，所以 $2as$ 是力。
